\section{Cournot Building Blocks}
\label{sec:product-market}

The private-adoption and government-incentive analysis needs only a small set of destination-level Cournot outcomes. We therefore collect the equilibrium blocks here and move the first-order-condition algebra, uniqueness checks, and interiority proofs to Appendix~\ref{app:cournot}. All expressions below are for the canonical baseline in which a firm that supports the local standard has zero marginal adaptation cost. The residual-cost extension is solved separately in Section~\ref{sec:selective-erosion}.

Let $P,A,B,C,D,H,S$ denote operating-profit blocks and $K_I,K_M,K_O,K_W$ the associated consumer-surplus blocks. Because markets are segmented, the same destination-level blocks can be assembled across formal regimes and adoption profiles without re-solving the full three-market game.

\subsection{Complete compatibility}
\label{subsec:product-market-complete}

When all three firms use the same standard at zero marginal cost, symmetry gives
\begin{equation}
    q_I=\frac{1}{4(1-v)},
    \label{eq:qI}
\end{equation}
and each firm earns
\begin{equation}
    P=\frac{1}{16(1-v)}.
    \label{eq:profit-P}
\end{equation}
Total quantity is $3q_I$, so consumer surplus is
\begin{equation}
    K_I=\frac{9}{32(1-v)^2}.
    \label{eq:KI}
\end{equation}

\subsection{Regional standardization}
\label{subsec:product-market-regional}

In a member market of a two-country regional standards union, the two member firms are mutually compatible and have zero marginal cost, while the outsider is a costly singleton. The equilibrium quantities are
\begin{equation}
    q_M=\frac{1+c}{2(2-3v)},
    \qquad
    q_B=\frac{1-3c-3(1-c)v}{2(2-3v)},
    \label{eq:quantities-su-member}
\end{equation}
with operating profits
\begin{equation}
    A=\frac{(1-v)(1+c)^2}{4(2-3v)^2},
    \qquad
    B=\frac{[1-3c-3(1-c)v]^2}{4(2-3v)^2}.
    \label{eq:profits-su-member}
\end{equation}
The corresponding consumer-surplus block is
\begin{equation}
    K_M
    =\frac{(3-c-3v+3cv)^2}{8(2-3v)^2}.
    \label{eq:KM}
\end{equation}

In the outsider market, the two coalition firms remain mutually compatible but each incurs marginal cost $c$, while the native outsider is a zero-cost singleton. The equilibrium quantities are
\begin{equation}
    q_C=\frac{1-2c}{2(2-3v)},
    \qquad
    q_D=\frac{1+2c-3v}{2(2-3v)},
    \label{eq:quantities-su-outsider}
\end{equation}
and the profit blocks are
\begin{equation}
    C=\frac{(1-v)(1-2c)^2}{4(2-3v)^2},
    \qquad
    D=\frac{(1+2c-3v)^2}{4(2-3v)^2}.
    \label{eq:profits-su-outsider}
\end{equation}
Consumer surplus is
\begin{equation}
    K_O
    =\frac{(3-2c-3v)^2}{8(2-3v)^2}.
    \label{eq:KO}
\end{equation}

\subsection{Separate national standards}
\label{subsec:product-market-sw}

Under separate national standards, all firms are compatibility singletons. In any destination the native firm has zero marginal cost and the two foreign firms have marginal cost $c$. Their quantities are
\begin{equation}
    q_H=\frac{1+2c}{4},
    \qquad
    q_S=\frac{1-2c}{4},
    \label{eq:quantities-sw}
\end{equation}
with operating profits
\begin{equation}
    H=\frac{(1+2c)^2}{16},
    \qquad
    S=\frac{(1-2c)^2}{16},
    \label{eq:profits-sw}
\end{equation}
and consumer surplus
\begin{equation}
    K_W=\frac{(3-2c)^2}{32}.
    \label{eq:KW}
\end{equation}

\subsection{Consumer surplus and regularity}
\label{subsec:product-market-blocks}

In every active-product configuration, the generalized price faced by the marginal consumer equals $1-X_k$, where $X_k$ is total quantity in market $k$. Hence
\begin{equation}
    CS_k
    =m_k\frac{X_k^2}{2}.
    \label{eq:consumer-surplus-general}
\end{equation}
This yields the four $K$ blocks above. Appendix~\ref{app:cournot} derives the underlying first-order systems and verifies uniqueness, positive quantities, and interior consumer participation on
\[
0<v<\frac14,
\qquad
0<c<\frac{1-3v}{3(1-v)}.
\]

\begin{lemma}[Canonical Cournot blocks]
\label{lem:cournot-blocks}
On the canonical parameter domain, the four product-market configurations used in the baseline continuation analysis have unique interior Cournot equilibria. Their quantities, operating profits, and consumer-surplus blocks are given by equations~\eqref{eq:qI}--\eqref{eq:KW}.
\end{lemma}

\input{tables/generated/table_cournot_blocks}

The economic distinction among the blocks is what matters for the remainder of the paper. $A$ and $B$ describe the protected member markets of a regional union, $C$ and $D$ describe the reciprocal outsider market, and $P$ is the common benchmark when all firms are compatible and costless. Section~\ref{sec:private-compatibility} uses these blocks to characterize private adoption; Section~\ref{sec:selective-erosion} then re-solves the member market when compatibility is complete but a positive residual adaptation cost remains.
